% Topic from _Linear Algebra_ Jim Hefferon
%  http://joshua.smcvt.edu/linearalgebra
%  2012-Feb-12
\topic{Pemeringkatan Halaman}
\index{pemeringkatan halaman|(}

Bayangkan Anda mencari buku terbaik tentang Aljabar Linear. Anda mungkin akan
mencobanya melalui mesin pencari web seperti Google. Mesin pencari menampilkan
halaman-halaman yang diurutkan menurut tingkat kepentingannya. Seperti dijelaskan
pendiri Google dalam \cite{BrinPage}, suatu halaman penting jika halaman penting
lain menaut kepadanya: ``suatu halaman dapat memiliki PageRank tinggi jika
banyak halaman menunjuk kepadanya, atau jika beberapa halaman yang menunjuk
kepadanya memiliki PageRank tinggi.'' Namun, bukankah definisi itu melingkar?
Bagaimana kepentingan suatu halaman dapat diketahui sebelum halaman-halaman
penting ditentukan? Jawabannya menggunakan nilai eigen dan vektor eigen.

Kita akan menyajikan versi sederhana algoritma PageRank. Untuk itu, kita
memodelkan World Wide Web sebagai kumpulan halaman yang dihubungkan oleh tautan.
Diagram dari \cite{Wills} berikut menggambarkan halaman sebagai lingkaran dan
tautan sebagai panah. Halaman~$p_1$ memiliki tautan ke halaman~$p_2$,
halaman~$p_2$ memiliki tautan ke~$p_3$, dan $p_3$ memiliki tautan ke $p_1$,
$p_2$, dan~$p_4$.
% \begin{center}  % add a little vertical spacing; looked tight to me.
%   \shortstack{\rule{0em}{1.5ex} \\ \includegraphics{jc/mp/ch5.9} \\ \rule{0em}{1.5ex}}
% \end{center}
\begin{center} 
  \includegraphics{jc/mp/ch5.9}
\end{center}

Gagasan utamanya adalah bahwa halaman patut diberi peringkat tinggi jika sering
dirujuk oleh halaman lain.
% (In practice people have tried to game the ranking system by
% setting up link farms of many pages that all point to each other.
% But our model will be simple.)
Artinya, kepentingan halaman~$p_i$ dinaikkan jika halaman~$p_j$ menaut kepadanya.
Besarnya kenaikan bergantung pada kepentingan halaman penaut~$p_j$, dibagi
dengan jumlah tautan keluar $a_j$ pada halaman tersebut.
\begin{equation*}
  \mathcal{I}(p_i)=\sum_{\text{halaman $p_j$ yang menaut masuk}}  \frac{\mathcal{I}(p_j)}{a_j}
\end{equation*}
Jadi, kepentingan $p_1$ sama dengan $1/3$ kali kepentingan~$p_3$ karena
satu-satunya tautan ke~$p_1$ berasal dari~$p_3$. Serupa dengan itu, kepentingan
$p_2$ adalah jumlah kepentingan~$p_1$ dan $1/3$ kali kepentingan~$p_3$.

Matriks berikut menyimpan informasi tersebut.
\begin{equation*}
  \begin{mat}
    0   &0  &1/3  &0   \\
    1   &0  &1/3  &0   \\
    0   &1  &0    &0 \\
    0   &0  &1/3  &0
  \end{mat}
\end{equation*}
Para pencipta algoritma itu memberikan cara berikut untuk memaknai matriks
tersebut.
\begin{quotation}
PageRank dapat dipandang sebagai model perilaku pengguna. Kita mengandaikan
adanya ``peselancar acak'' yang diberi suatu halaman web secara acak dan terus
mengeklik tautan tanpa pernah menekan ``kembali'' \ldots{}
% but eventually gets bored
% and starts on another random page. 
Peluang peselancar acak mengunjungi suatu halaman adalah PageRank halaman itu.
\cite{BrinPage}
\end{quotation}
Jadi, dari baris pertama matriks terlihat bahwa satu-satunya cara peselancar
acak mencapai $p_1$ adalah datang dari~$p_3$. Halaman itu memiliki tiga tautan,
sehingga peluang mengeklik tautan menuju $p_1$ adalah $1/3$ kali peluang berada
di halaman~$p_3$.

Hal ini memunculkan pertanyaan tentang halaman~$p_4$. Di Internet, banyak
halaman merupakan \definend{tautan menggantung}%
\index{tautan!menggantung}\index{tautan menggantung} atau
\definend{tautan penyerap},\index{tautan!penyerap}\index{tautan penyerap}
tanpa tautan keluar. Apa yang terjadi pada peselancar acak yang mengunjungi
halaman demikian? Cara paling sederhana adalah membayangkan bahwa peselancar
memilih halaman berikutnya sepenuhnya secara acak.
\begin{equation*}
  H=\begin{mat}
    0   &0  &1/3  &1/4   \\
    1   &0  &1/3  &1/4   \\
    0   &1  &0    &1/4 \\
    0   &0  &1/3  &1/4
  \end{mat}
\end{equation*}

Kita mencari vektor $\vec{\mathcal{I}}$ yang komponennya merupakan peringkat
kepentingan setiap halaman $\mathcal{I}(p_i)$. Dengan notasi ini, syarat
pemeringkatan halaman adalah $H\vec{\mathcal{I}}=\vec{\mathcal{I}}$. Artinya,
kita menginginkan vektor eigen matriks yang terkait dengan nilai
eigen~$\lambda=1$.

Berikut perhitungan vektor eigen oleh \textit{Sage} (disunting agar muat pada
halaman).
\begin{lstlisting}
sage: H=matrix([[0,0,1/3,1/4], [1,0,1/3,1/4], [0,1,0,1/4], [0,0,1/3,1/4]])   
sage: H.eigenvectors_right()                                                 
[(1, [(1, 2, 9/4, 1)], 1), 
 (0, [(0, 1, 3, -4)], 1), 
 (-0.3750000000000000? - 0.4389855730355308?*I, 
     [(1, -0.1250000000000000? + 1.316956719106593?*I, 
          -1.875000000000000? - 1.316956719106593?*I, 1)], 1), 
 (-0.3750000000000000? + 0.4389855730355308?*I, 
     [(1, -0.1250000000000000? - 1.316956719106593?*I, 
          -1.875000000000000? + 1.316956719106593?*I, 1)], 1)]
\end{lstlisting}
Vektor eigen yang diberikan \textit{Sage} untuk nilai eigen~$\lambda=1$ adalah
sebagai berikut.
\begin{equation*}
  \colvec{1 \\ 2 \\ 9/4 \\ 1}
\end{equation*}
Tentu terdapat banyak vektor dalam ruang eigen tersebut. Karena komponen
PageRank adalah peluang, kita normalkan agar jumlah komponennya satu.
\begin{lstlisting}
sage: v=vector([1, 2, 9/4, 1])
sage: v/sum(v)
(4/25, 8/25, 9/25, 4/25)
sage: w=v/sum(v)
sage: w.n()
(0.160000000000000, 0.320000000000000, 0.360000000000000, 0.160000000000000)
\end{lstlisting}
Jadi, halaman pertama dan keempat diberi tingkat kepentingan yang sama. Halaman
kedua dan ketiga jauh lebih penting daripada keduanya dan kira-kira sama penting
satu sama lain.

Kita tambahkan satu penyempurnaan lagi. Peselancar diperbolehkan memilih halaman
baru secara acak meskipun tidak sedang berada di halaman menggantung. Dalam
persamaan berikut, hal itu terjadi dengan peluang~$1-\alpha$.
\begin{equation*}
  G=\alpha\cdot\begin{mat}
    0   &0  &1/3  &1/4   \\
    1   &0  &1/3  &1/4   \\
    0   &1  &0    &1/4 \\
    0   &0  &1/3  &1/4
  \end{mat}
  +(1-\alpha)\cdot\begin{mat}
    1/4   &1/4  &1/4  &1/4   \\
    1/4   &1/4  &1/4  &1/4   \\
    1/4   &1/4  &1/4  &1/4  \\
    1/4   &1/4  &1/4  &1/4
  \end{mat}
\end{equation*}
Ini disebut \definend{matriks Google}\index{matriks Google}%
\index{matriks!Google}.

Dalam praktik, $\alpha$ biasanya berada di antara $0.85$ dan~$0.99$. Berikut
peringkat keempat halaman untuk beberapa nilai $\alpha$.
\begin{center}
  \begin{tabular}{r|cccc}
    \multicolumn{1}{c}{$\alpha$}
          &$0.85$  &$0.90$  &$0.95$ &$0.99$ \\ \hline
    $p_1$ &$0.171$ &$0.167$ &$0.163$  &$0.161$  \\
    $p_2$ &$0.316$ &$0.317$ &$0.319$  &$0.320$  \\
    $p_3$ &$0.342$ &$0.348$ &$0.354$  &$0.359$  \\
    $p_4$ &$0.171$ &$0.167$ &$0.163$  &$0.161$  
  \end{tabular}
\end{center}
% Numbers computed as:
% sage: H=matrix(QQ,[[0,0,1/3,1/4], [1,0,1/3,1/4], [0,1,0,1/4], [0,0,1/3,1/4]])
% sage: S=matrix(QQ,[[1/4,1/4,1/4,1/4], [1/4,1/4,1/4,1/4], [1/4,1/4,1/4,1/4], [1/4,1/4,1/4,1/4]]) 
% sage: I=matrix(QQ,[[1,0,0,0], [0,1,0,0], [0,0,1,0], [0,0,0,1]])
% sage: alpha=0.95                       
% sage: G=alpha*H+(1-alpha)*S    
% sage: N=G-I  
% sage: 1200*N         
% [-1185.00000000000  15.0000000000000  395.000000000000  300.000000000000]
% [ 1155.00000000000 -1185.00000000000  395.000000000000  300.000000000000]
% [ 15.0000000000000  1155.00000000000 -1185.00000000000  300.000000000000]
% [ 15.0000000000000  15.0000000000000  395.000000000000 -900.000000000000]
% sage: M=matrix(QQ,[[-1185,15,395,300], [1155,-1185,395,300], [15,1155,-1185,300], [15,15,395,-900]])
% sage: M.echelon_form()         
% [         1          0          0         -1]
% [         0          1          0     -39/20]
% [         0          0          1 -3423/1580]
% [         0          0          0          0]
% sage: v=vector([1,39/20,3423/1580,1]) 
% sage: v/v.norm().n()   
% (0.308665053665036, 0.601896854646821, 0.668709163731278, 0.308665053665036)


\medskip
Rincian algoritma yang digunakan mesin pencari komersial bersifat rahasia,
memiliki banyak penyempurnaan, dan sering berubah. Namun, para pencipta Google
berkenan menguraikan dasar pekerjaan mereka dalam \cite{BrinPage}. Sumber yang
lebih mutakhir adalah \cite{WikipediaPageRank}. Dua paparan bagus lainnya adalah
\cite{Wills} dan \cite{Austin}.

\begin{exercises}
  \item Matriks persegi disebut \definend{stokastik}\index{matriks stokastik}%
    \index{matriks!stokastik} jika jumlah entri pada setiap kolom adalah satu.
    Matriks Google dihitung sebagai kombinasi
    $G=\alpha*H+(1-\alpha)*S$ dari dua matriks stokastik.
    Tunjukkan bahwa $G$ harus stokastik.
    \begin{answer}
      Jumlah entri pada kolom~$j$ adalah
      $\sum_i \alpha h_{i,j}+(1-\alpha)s_{i,j}
       =\sum_i \alpha h_{i,j}+\sum_i (1-\alpha)s_{i,j}
       =\alpha\sum_i h_{i,j} +(1-\alpha)\sum_i s_{i,j}
       =\alpha\cdot 1+(1-\alpha)\cdot 1$,
      yaitu satu.
    \end{answer}
  \item Untuk jejaring halaman berikut, kepentingan setiap halaman seharusnya
    sama. Verifikasikan hal itu untuk $\alpha=0.85$.
    \begin{center}
      \includegraphics{jc/mp/ch5.11}
    \end{center}
    \begin{answer}
      Sesi \textit{Sage} berikut memberikan nilai-nilai yang sama.
\begin{lstlisting}
sage: H=matrix(QQ,[[0,0,0,1], [1,0,0,0], [0,1,0,0], [0,0,1,0]])
sage: S=matrix(QQ,[[1/4,1/4,1/4,1/4], [1/4,1/4,1/4,1/4], [1/4,1/4,1/4,1/4], [1/4,1/4,1/4,1/4]])
sage: alpha=0.85     
sage: G=alpha*H+(1-alpha)*S    
sage: I=matrix(QQ,[[1,0,0,0], [0,1,0,0], [0,0,1,0], [0,0,0,1]])
sage: N=G-I
sage: 1200*N         
[-1155.00000000000  45.0000000000000  45.0000000000000  1065.00000000000]
[ 1065.00000000000 -1155.00000000000  45.0000000000000  45.0000000000000]
[ 45.0000000000000  1065.00000000000 -1155.00000000000  45.0000000000000]
[ 45.0000000000000  45.0000000000000  1065.00000000000 -1155.00000000000]
sage: M=matrix(QQ,[[-1155,45,45,1065], [1065,-1155,45,45], [45,1065,-1155,45], [45,45,1065,-1155]]) 
sage: M.echelon_form()       
[ 1  0  0 -1]
[ 0  1  0 -1]
[ 0  0  1 -1]
[ 0  0  0  0]
sage: v=vector([1,1,1,1])
sage: (v/v.norm()).n()
(0.500000000000000, 0.500000000000000, 0.500000000000000, 0.500000000000000)
\end{lstlisting}
    \end{answer}
  \item \cite{BryanLeise}
    Berikan peringkat kepentingan untuk jejaring halaman berikut.
    \begin{center}
      \includegraphics{jc/mp/ch5.10}
    \end{center}
    \begin{exparts}
      \item Gunakan $\alpha=0.85$.
      \item Gunakan $\alpha=0.95$.
      \item Perhatikan bahwa walaupun semua halaman lain menaut ke $p_3$,
        sehingga halaman itu tampak penting, $p_3$ bukan halaman berperingkat
        tertinggi. Halaman manakah yang berperingkat tertinggi? Jelaskan.
    \end{exparts}
    \begin{answer}
      Kita mempunyai matriks berikut.
      \begin{equation*}
          H=\begin{mat}
               0     &0    &1    &1/2   \\
               1/3   &0    &0    &0   \\
               1/3   &1/2  &0    &1/2 \\
               1/3   &1/2  &0    &0
          \end{mat}
      \end{equation*}
      \begin{exparts}
        \item Sesi \Sage{} berikut memberikan jawabannya.
\begin{lstlisting}
sage: H=matrix(QQ,[[0,0,1,1/2], [1/3,0,0,0], [1/3,1/2,0,1/2], [1/3,1/2,0,0]])
sage: S=matrix(QQ,[[1/4,1/4,1/4,1/4], [1/4,1/4,1/4,1/4], [1/4,1/4,1/4,1/4], [1/4,1/4,1/4,1/4]])
sage: I=matrix(QQ,[[1,0,0,0], [0,1,0,0], [0,0,1,0], [0,0,0,1]])
sage: alpha=0.85
sage: G=alpha*H+(1-alpha)*S
sage: N=G-I
sage: 1200*N
[-1155.00000000000  45.0000000000000  1065.00000000000  555.000000000000]
[ 385.000000000000 -1155.00000000000  45.0000000000000  45.0000000000000]
[ 385.000000000000  555.000000000000 -1155.00000000000  555.000000000000]
[ 385.000000000000  555.000000000000  45.0000000000000 -1155.00000000000]
sage: M=matrix(QQ,[[-1155,45,1065,555], [385,-1155,45,45], [385,555,-1155,555],
....:      [385,555,45,-1155]])
sage: M.echelon_form()
[            1             0             0 -106613/58520]
[            0             1             0        -40/57]
[            0             0             1        -57/40]
[            0             0             0             0]
sage: v=vector([106613/58520,40/57,57/40,1])
sage: (v/sum(v)).n()
(0.368150677047603, 0.141809358496821, 0.287961628597607, 0.202078335857970)
\end{lstlisting}
        \item Lanjutkan sesi \Sage{} untuk memperoleh hasil berikut.
\begin{lstlisting}
 sage: alpha=0.95
sage: G=alpha*H+(1-alpha)*S
sage: N=G-I
sage: 1200*N
[-1185.00000000000  15.0000000000000  1155.00000000000  585.000000000000]
[ 395.000000000000 -1185.00000000000  15.0000000000000  15.0000000000000]
[ 395.000000000000  585.000000000000 -1185.00000000000  585.000000000000]
[ 395.000000000000  585.000000000000  15.0000000000000 -1185.00000000000]
sage: M=matrix(QQ,[[-1185,15,1155,585], [395,-1185,15,15], [395,585,-1185,585],
....:      [395,585,15,-1185]])
sage: M.echelon_form()
[             1              0              0 -361677/186440]
[             0              1              0         -40/59]
[             0              0              1         -59/40]
[             0              0              0              0]
sage: v=vector([361677/186440,40/59,59/40,1])
sage: (v/sum(v)).n()
(0.380906693515433, 0.133120452946554, 0.289620185441846, 0.196352668096167)
\end{lstlisting}
       \item Halaman $p_3$ penting, tetapi meneruskan seluruh kepentingannya
         hanya kepada satu halaman, yaitu $p_1$. Karena itu, halaman tersebut
         menerima kenaikan besar.
      \end{exparts}    
    \end{answer}
\end{exercises}
\index{pemeringkatan halaman|)}
\endinput
